\label{annexe:ingestion}

Ce script Python est responsable de la transformation des données non structurées en Graphe de Connaissances. Il utilise les modèles \texttt{Gemini/Llama 3} pour extraire les entités et relations.

\begin{lstlisting}[language=python, caption=Script d'ingestion et de construction du Graphe (ingest\_graph.py)]
# Configuration Neo4j
os.environ["NEO4J_URI"] = "bolt://neo4j:7687"
os.environ["NEO4J_USERNAME"] = "neo4j"
os.environ["NEO4J_PASSWORD"] = "password1234"

# Donnees brutes 
raw_text = [
    """
    """
]

def ingest_graph():
    print("Connexion a Neo4j...")
    graph = Neo4jGraph()

    provider = os.getenv("LLM_PROVIDER", "google").lower()
    if provider == "groq":
        print("Ingestion via GROQ...")
        llm = ChatGroq(
            model=os.getenv("GROQ_MODEL", "llama-3.3-70b-versatile"),
            temperature=0
        )
    else:
        print("Ingestion via GOOGLE...")
        llm = ChatGoogleGenerativeAI(
            model=os.getenv("GOOGLE_MODEL", "gemini-2.5-flash"),
            temperature=0
        )
    
    llm_transformer = LLMGraphTransformer(llm=llm)

    print("Analyse du texte (Gemini)...")
    documents = [Document(page_content=text) for text in raw_text]
    
    # Transformation
    try:
        graph_documents = llm_transformer.convert_to_graph_documents(documents)
    except Exception as e:
        print(f"Erreur Extraction LLM: {e}")
        return

    # --- AFFICHAGE DU RESUME AVANT INSERTION ---
    nodes = graph_documents[0].nodes
    rels = graph_documents[0].relationships
    
    print(f"\n RESUME DE L'EXTRACTION :")
    print(f"   Noeuds trouves : {len(nodes)}")
    print(f"   Relations trouvees : {len(rels)}")
    
    print("\n APERÇU DES RELATIONS (5 premieres) :")
    for i, r in enumerate(rels[:5]):
        print(f"   ({r.source.id}) --[{r.type}]--> ({r.target.id})")

    print("\n Sauvegarde dans Neo4j...")
    graph.add_graph_documents(graph_documents)
    graph.refresh_schema()
    
    count_query = "MATCH (n) RETURN count(n) as count"
    result = graph.query(count_query)
    print(f"   Total Nœuds en base : {result[0]['count']}")
    
    print("\n Ingestion terminee avec succes !")

if __name__ == "__main__":
    ingest_graph()
\end{lstlisting}
